%----------------------------------------------------------------------
\section{Preliminaries}\label{sec:prelim}
%----------------------------------------------------------------------

\subsection{Notation inherited from [DAF]}

Throughout, $V = \GF(2)^n$ denotes the discriminator space;
$\bigwedge V$ the exterior algebra; $\kappa$ a (finite,
ordered) discrimination regime; $\Omega = \GF(2)^2$ the
subobject classifier. We write $f[S] \in \GF(2)$ for the
coefficient of basis element $e_S$ in an exterior element
$f \in \bigwedge V$, indexing subsets $S \subseteq [n]$ by
their $n$-bit characteristic masks. $\oplus$ is GF(2)
addition (XOR); $\wedge$ is the exterior wedge;
$\partial: \bigwedge^k V \to \bigwedge^{k-1} V$ is the
boundary operator of [DAF] \S5.

\subsection{The encoding lattice}
\label{sec:encoding}

Any $f \in \bigwedge V$ has two linearly equivalent dense
representations and several compressed variants.

\begin{definition}[Coefficient basis $\Lambda$]\label{def:lambda}
  The \emph{coefficient representation} of $f \in \bigwedge V$
  is the vector $(f[S])_{S \subseteq [n]} \in \GF(2)^{2^n}$.
  Multiplication by a basis element is a single-coordinate
  update; wedge product is disjoint-pair convolution;
  boundary is gather--XOR over supersets.
\end{definition}

\begin{definition}[Zeta transform]\label{def:zeta}
  The \emph{zeta transform} $\zeta: \GF(2)^{2^n} \to
  \GF(2)^{2^n}$ is defined on a function $f$ by
  \begin{equation}
    \zeta(f)[T] = \bigoplus_{S \subseteq T} f[S],
    \qquad T \subseteq [n].
  \end{equation}
\end{definition}

\begin{proposition}[Zeta involutivity over $\GF(2)$]
  \label{prop:zeta-involution}
  \footnote{Class \textsf{A}. Warrant: characteristic 2
  collapses the M\"{o}bius coefficient $(-1)^{|T| - |S|}$
  to the constant 1, so the M\"{o}bius inverse and the
  zeta transform coincide. Qualifier: over fields of
  characteristic $\neq 2$ the two operations are distinct.
  Finite-$\kappa$ valid.}
  Over $\GF(2)$, $\zeta \circ \zeta = \mathrm{id}$.
\end{proposition}

\begin{proof}
  $\zeta^2(f)[T] = \bigoplus_{S \subseteq T}
   \bigoplus_{R \subseteq S} f[R] =
   \bigoplus_{R \subseteq T} N(R,T) \cdot f[R]$
  where $N(R,T) = |\{S : R \subseteq S \subseteq T\}|
  = 2^{|T| - |R|}$. Over $\GF(2)$, $N(R,T) \equiv 1$ iff
  $|T| = |R|$, i.e.\ iff $R = T$. Hence
  $\zeta^2(f)[T] = f[T]$.
\end{proof}

\begin{definition}[Reed--Muller syndrome basis]\label{def:rm}
  The \emph{RM representation} of $f$ is $\zeta(f) \in
  \GF(2)^{2^n}$. By Proposition \ref{prop:zeta-involution},
  the inverse transform is $\zeta$ itself.
\end{definition}

\begin{definition}[Grade-packed basis]\label{def:grade-packed}
  For $r \in \{0, 1, \ldots, n\}$, the
  \emph{grade-$\leq r$-packed basis} represents $f$ by
  its grade-$\leq r$ coefficients only:
  $(f[S])_{|S| \leq r}$, of length
  $R_r = \sum_{k=0}^{r} \binom{n}{k}$. It is lossy on
  general $f$ and lossless on $f$ whose coefficients
  vanish above grade $r$.
\end{definition}

\begin{definition}[Hamming-compressed basis]\label{def:hamming}
  At $n = 3$: $\bigwedge V \setminus \{0\}$ has 7 nonzero
  basis elements, corresponding to the 7 columns of the
  $[7,4,3]$ Hamming parity-check matrix $H$. An element
  $f$ whose nonzero support lies in a coset of the
  Hamming information subspace is represented by its 4
  information bits plus an implicit 3-bit syndrome.
  The generalization to $n > 3$ is Reed--Muller
  $\mathrm{RM}(r, n)$ of the chosen grade-$r$ subcode.
\end{definition}

We collect these four representations into an ordered
lattice of encodings $\mathcal{E} = \{\Lambda,
\Lambda^{\leq r}, \mathrm{RM}, \mathrm{RM}^{\leq r},
\mathrm{Ham}\}$, with the partial order generated by the
observation that grade-truncated forms are projections
of their full-grade parents.

\begin{proposition}[Encoding--addition commutation]
  \label{prop:encoding-commute}
  \footnote{Class \textsf{A}. Warrant: $\zeta$ is
  $\GF(2)$-linear; addition in $\bigwedge V$ is
  pointwise XOR; linearity commutes with pointwise XOR.
  Finite-$\kappa$ valid.}
  For $f, g \in \bigwedge V$,
  $\zeta(f \oplus g) = \zeta(f) \oplus \zeta(g)$.
  Hence $\oplus$ has a single computational form that
  is correct in both $\Lambda$ and $\mathrm{RM}$.
\end{proposition}

\begin{proof}
  Direct from $\GF(2)$-linearity of summation over
  subsets.
\end{proof}

\subsection{The resource vector}\label{sec:resource-vector}

Every production in the grammar of Section
\ref{sec:grammar} carries a cost attribute
$c \in \mathbb{N}^4$ ordered as
$c = (\VRAM, \SMEM, \REG, \LAT)$.

\begin{definition}[Resource vector]\label{def:cost-vector}
  The \emph{resource vector} of a plan $P$ is
  $c(P) = (\VRAM(P), \SMEM(P), \REG(P), \LAT(P))
  \in \mathbb{N}^4$, where
  \begin{description}[nosep,leftmargin=0.9cm]
    \item[$\VRAM(P)$] peak device memory used by $P$,
      in bytes, including inputs, outputs, and any
      intermediates materialized outside the SM;
    \item[$\SMEM(P)$] peak per-SM shared-memory
      occupancy, in bytes, across any single kernel
      that $P$ dispatches;
    \item[$\REG(P)$] peak per-thread register live-set
      size, in words;
    \item[$\LAT(P)$] expected wall-clock latency, in
      any chosen time unit consistent across the
      hardware profile.
  \end{description}
\end{definition}

\begin{definition}[Budget vector]\label{def:budget}
  A \emph{budget vector} is a vector $B \in
  \mathbb{N}^4$ in the same order as
  Definition \ref{def:cost-vector}. A plan $P$ is
  \emph{admissible under $B$} if $c(P)_r \leq B_r$
  for each resource $r \in \{\VRAM, \SMEM, \REG,
  \LAT\}$.
\end{definition}

\begin{definition}[Cost lattice]\label{def:cost-lattice}
  Cost vectors form a join-semilattice under the
  componentwise partial order $c \leq c'$ iff $c_r \leq
  c'_r$ for all $r$. The join $c \vee c'$ is the
  componentwise max. We refer to the antichain of
  Pareto-minimal elements of a set of cost vectors as
  the \emph{Pareto frontier}.
\end{definition}

Cost combination under the grammar's productions uses
different operations per dimension, depending on
whether the production is sequential (\LAT\ adds,
\REG\ sometimes maxes) or parallel (\LAT\ maxes,
\REG\ sums). Section \ref{sec:grammar} states the
exact rule per production.

\subsection{Hardware profile}\label{sec:profile}

\begin{definition}[Hardware profile]\label{def:profile}
  A \emph{hardware profile} is a function $\pi: \mathcal{A}
  \times \mathbb{N} \to \mathbb{N}^4$ assigning a
  resource vector to each atom $a \in \mathcal{A}$
  at each batch size $k$. A profile is \emph{monotone}
  if each component is non-decreasing in batch size,
  and \emph{sub-additive in latency} if the latency of
  concurrent dispatch of independent atoms is at most
  the max of individual latencies plus a constant
  launch overhead. All theorems of this paper assume
  a monotone, sub-additive profile.
\end{definition}

A concrete profile for a reference GPU is deferred to
follow-up work (\S1 item S1). The grammar and parser
are architecture-independent; only the profile changes
between devices.

\subsection{Tile geometry and LCM packing}

\begin{definition}[Tile]\label{def:tile}
  A \emph{tile} of size $d$ is a $d \times d$ block of
  $\GF(2)$ values. The \emph{natural tile size} of an
  atom $a$ is the smallest $d_a$ such that $a$'s input
  and output fit in a single $d_a \times d_a$ block;
  e.g.\ $d_{\mathrm{rotate}} = 2$,
  $d_{\mathrm{fano}} = 3$,
  $d_{\mathrm{boundary}_n} = n$.
\end{definition}

\begin{definition}[Hardware tile]\label{def:hw-tile}
  A \emph{hardware tile} of size $M$ is a
  $M \times M$ block that the target device treats as
  a primitive dispatch unit (e.g.\ an NVIDIA WMMA tile
  at $M = 16$, or a warp-level sub-tile at $M = 4$).
\end{definition}

\begin{proposition}[LCM packing density]
  \label{prop:lcm-density}
  \footnote{Class \textsf{AH}. Warrant: arithmetic
  identity. Qualifier: assumes the hardware profile
  permits stride-$d$ access within an $M \times M$
  block at no bandwidth cost, as is the case for
  contemporary tensor-core pipelines over $\GF(2)$-emulating
  integer arithmetic. Finite-$\kappa$ valid.}
  Let $d$ be a natural tile size, $M$ a hardware tile
  size, and $L = \lcm(d, M)$. In an $L \times L$
  block, there are $(L/d)^2$ CSTZ tiles and $(L/M)^2$
  hardware tiles. Each hardware tile holds $(M/d)^2$
  CSTZ tiles (integer when $d \mid M$, rational
  otherwise in a sense made precise below), and the
  information density of a hardware tile is
  exactly $M^2$ bits, independent of $d$.
\end{proposition}

\begin{proof}
  Total CSTZ bits in the $L \times L$ block:
  $(L/d)^2 \cdot d^2 = L^2$. Total hardware tiles:
  $(L/M)^2$. Bits per hardware tile: $L^2 / (L/M)^2
  = M^2$. The per-tile CSTZ op count is
  $((L/d)^2) / ((L/M)^2) = (M/d)^2$ in the rational
  sense; when $d \mid M$ this is an integer, when
  $d \nmid M$ it is the average over a block that has
  integer counts that fluctuate by $\pm 1$.
\end{proof}

\begin{remark}[Forced batch coarseness]
  \label{rem:coarseness}
  Proposition \ref{prop:lcm-density} says that if the
  planner is willing to submit a batch of size at
  least $(L/d)^2$ CSTZ ops, it saturates hardware
  tiles at 100\% density regardless of $d$. Below
  $(L/d)^2$, dispatch is sparse and some hardware
  lanes carry no CSTZ work; Section \ref{sec:portfolios}
  shows how the planner repurposes those empty lanes
  via portfolio packing.
\end{remark}

\subsection{Terminology concordance}
\label{sec:concordance}

The paper inherits CSTZ terminology from [DAF] and
draws compiler-theoretic terminology from several
established frameworks.

\emph{Compiler theory.}
\begin{description}[nosep,leftmargin=1.5cm]
  \item[Atom] A terminal of the grammar, representing
    one kernel-level primitive. This is the compiler's
    analogue of a \emph{machine instruction} in
    RISC ISA design (Hennessy \& Patterson 2019).
  \item[Production] A rewrite rule combining atoms or
    subplans into a larger plan. Analogue of a
    \emph{scheduling transformation} in
    polyhedral-model compilers.
  \item[Plan] A derivation of the workload's output
    nonterminal from atoms, corresponding to a
    specific executable schedule. Analogue of the
    compiler's \emph{intermediate representation}
    after lowering.
\end{description}

\emph{Parser theory} (Earley 1970; Aycock \& Horspool 2002).
\begin{description}[nosep,leftmargin=1.5cm]
  \item[Chart item] A record $(A \to \alpha \bullet
    \beta, i, j, \bar{c})$ where $(A \to \alpha \beta)$
    is a production, the dot separates matched from
    unmatched children, $[i, j]$ is the range of
    workload atoms consumed, and $\bar{c}$ is the
    Pareto-minimal cost vector achieving that dotted
    item. Standard Earley augmented with an attribute
    in the style of Knuth (1968).
  \item[Completer] The Earley step combining a
    finished item $(A \to \alpha \bullet, i, j, \bar{c})$
    with a waiting item $(B \to \beta \bullet A \gamma,
    h, i, \bar{c}')$ to produce
    $(B \to \beta A \bullet \gamma, h, j, \bar{c}
    \sqcup \bar{c}')$ where $\sqcup$ is the
    production-appropriate cost combination.
  \item[Admissible heuristic] A function estimating a
    lower bound on the cost-to-complete any subplan.
    Must be \emph{monotone} (the triangle inequality)
    for A$^*$ correctness (Dechter \& Pearl 1985).
\end{description}

\emph{GPU architecture.}
\begin{description}[nosep,leftmargin=1.5cm]
  \item[Warp / lane] A group of (typically 32) lanes
    executing in lockstep under SIMT. We use the
    stronger MIMD assumption (Section \ref{sec:fusion}),
    available on current hardware via masked execution:
    lanes all execute the same instruction, but the
    mask encodes different masked-matrix entries per
    lane, giving lane-local behavior that MIMD would
    provide.
  \item[SM] A streaming multiprocessor --- the physical
    unit holding one shared-memory region and one
    register file. Per-SM resource budgets (shared
    memory, register file) apply to any single kernel
    dispatched to that SM.
  \item[Kernel] A function compiled for the GPU, whose
    invocation is the atomic unit of the three-tier
    memory budget. A fused kernel is one kernel
    dispatch that performs work spanning multiple
    CSTZ atoms.
\end{description}

\emph{Operational semantics.} We adopt a
step-indexed operational semantics for plans,
whose configurations are triples $(W, \chi, s)$
with $W$ the remaining workload, $\chi$ the chart,
and $s$ the register state. Transitions are parser
operations (predict, scan, complete). This is
entirely standard (Plotkin 1981).
